\ifdefined\gkver\else\def\gkver{student}\fi
\documentclass[\gkver]{gaokaozhenti}
\begin{document}

\chapter{2008年广东理综}
%% paperType: 理综物理部分

\begin{enumerate}
\item
%% number: 1
%% paperName: 2008年广东理综第 1 题 4 分
%% typeId: 1
%% type: 单选题
%% chapter: 直线运动
%% point: 匀变速直线运动
%% method: 无
%% score: 4
%% degree: 850
%% duplicateId: 0
%% body:
伽利略在著名的斜面实验中，让小球分别从倾角不同、阻力很小的斜面从静止开始滚下，他通过实验观察和逻辑推理，得到的正确结论有\xzanswer{B}

\fourchoices[answer=B]
{倾角一定时，小球在斜面上的位移与时间成正比}
{倾角一定时，小球在斜面上的速度与时间成正比}
{斜面长一定时，小球从顶端滚到底端时的速度与倾角无关}
{斜面长一定时，小球从顶端滚到底端所需的时间与倾角无关}

%% answer: B
%% memo:
\memoanswer{倾角一定时，小球在斜面上的位移与时间的平方成正比，在斜面上的速度与时间成正比，故选项 A 错误，选项 B 正确。斜面长度一定时，小球从顶端滚到底端时的速度与倾角有关，从顶端滚到底端所需时间与倾角有关，故选项 C、D 错误。}

\item
%% number: 2
%% paperName: 2008年广东理综第 2 题 4 分
%% typeId: 2
%% type: 多选题
%% chapter: 原子和原子核
%% point: 人工核反应
%% method: 无
%% score: 4
%% degree: 850
%% duplicateId: 0
%% body:
铝箔被 $\alpha$ 粒子轰击后发生了以下核反应：$\ce{^{27}_{13}Al + ^{4}_{2}He -> X + ^{1}_{0}n}$。下列判断正确的是\xzanswer{BD}

\fourchoices[answer=BD]
{$\ce{^{1}_{0}n}$ 是质子}
{$\ce{^{1}_{0}n}$ 是中子}
{X 是 $\ce{^{28}_{14}Si}$ 的同位素}
{X 是 $\ce{^{31}_{15}P}$ 的同位素}

%% answer: BD
%% memo:
\memoanswer{根据核反应方程质量数和电荷数守恒可知选项 B、D 正确。}

\item
%% number: 3
%% paperName: 2008年广东理综第 3 题 4 分
%% typeId: 1
%% type: 单选题
%% chapter: 功和能
%% point: 动能定理
%% method: 无
%% score: 4
%% degree: 850
%% duplicateId: 0
%% body:
运动员跳伞将经历加速下降和减速下降两个过程。将人和伞看成一个系统，在这两个过程中，下列说法正确的是\xzanswer{A}

\fourchoices[answer=A]
{阻力对系统始终做负功}
{系统受到的合外力始终向下}
{重力做功使系统的重力势能增加}
{任意相等的时间内重力做的功相等}

%% answer: A
%% memo:
\memoanswer{在两个过程中，阻力始终对系统做负功，选项 A 正确。加速下降时，系统受到的合力向下；减速下降时，系统受到的合力向上，选项 B 错误。两个过程中，重力始终做正功，系统的重力势能减少，选项 C 错误。在任意相等时间内，系统下降的高度不相等，故重力做功不相等，选项 D 错误。}

\item
%% number: 4
%% paperName: 2008年广东理综第 4 题 4 分
%% typeId: 2
%% type: 多选题
%% chapter: 磁场
%% point: 带电粒子在磁场中的运动
%% method: 无
%% score: 4
%% degree: 650
%% duplicateId: 0
%% body:
1930 年劳伦斯制成了世界上第一台回旋加速器。其原理如图所示，这台加速器由两个铜质 D 形盒 $D_{1}$、$D_{2}$ 构成，其间留有空隙。下列说法正确的是\xzanswer{AD}

\onepicture[w=4.0cm]{figs/04.png}

\fourchoices[answer=AD]
{离子由加速器的中心附近进入加速器}
{离子由加速器的边缘进入加速器}
{离子从磁场中获得能量}
{离子从电场中获得能量}

%% answer: AD
%% memo:
\memoanswer{离子由加速器的中心附近进入加速器，从电场中获取能量，最后从加速器边缘离开加速器，选项 A、D 正确。}

\item
%% number: 5
%% paperName: 2008年广东理综第 5 题 4 分
%% typeId: 1
%% type: 单选题
%% chapter: 交流电
%% point: 交流电
%% method: 无
%% score: 4
%% degree: 650
%% duplicateId: 0
%% body:
小型交流发电机中，矩形金属线框在匀强磁场中匀速转动，产生的感应电动势与时间呈正弦函数关系，如图所示，此线圈与一个 $R=10\UO$ 的电阻构成闭合回路，不计电路的其他电阻，下列说法正确的是\xzanswer{D}

\onepicture[w=5.0cm]{figs/05.png}

\fourchoices[answer=D]
{交变电流的周期为 $0.125\Us$}
{交变电流的频率为 $8\UHz$}
{交变电流的有效值为 $\sqrt{2}\UA$}
{交变电流的最大值为 $4\UA$}

%% answer: D
%% memo:
\memoanswer{由 $e-t$ 图像可知，交变电流的周期为 $0.25\Us$，故频率为 $4\UHz$，选项 A、B 错误。根据欧姆定律可知交变电流的最大值为 $2\UA$，故有效值为 $\sqrt{2}\UA$，选项 C 正确。（本题题后解析按图像所得结论为选项 C，而原卷参考答案页与平台数据均标 D，二者不一致；答案命令按平台数据取 D，解析照原卷如实保留。）}

\item
%% number: 6
%% paperName: 2008年广东理综第 6 题 4 分
%% typeId: 2
%% type: 多选题
%% chapter: 原子和原子核
%% point: 玻尔模型
%% method: 无
%% score: 4
%% degree: 650
%% duplicateId: 0
%% body:
有关氢原子光谱的说法正确的是\xzanswer{BC}

\fourchoices[answer=BC]
{氢原子的发射光谱是连续谱}
{氢原子光谱说明氢原子只发出特定频率的光}
{氢原子光谱说明氢原子能级是分立的}
{氢原子光谱线的频率与氢原子能级的能量差无关}

%% answer: BC
%% memo:
\memoanswer{氢原子的发射光谱是不连续的，只能发出特定频率的光，说明氢原子的能级是分立的，选项 B、C 正确。根据玻尔理论可知，选项 D 错误。}

\item
%% number: 7
%% paperName: 2008年广东理综第 7 题 4 分
%% typeId: 1
%% type: 单选题
%% chapter: 电路
%% point: 动态分析
%% method: 无
%% score: 4
%% degree: 650
%% duplicateId: 0
%% body:
电动势为 $E$、内阻为 $r$ 的电源与定值电阻 $R_{1}$、$R_{2}$ 及滑动变阻器 $R$ 构成如图所示的电路。当滑动变阻器的触头由中点滑向 b 端时，下列说法正确的是\xzanswer{A}

\onepicture[w=5.0cm]{figs/07.png}

\fourchoices[answer=A]
{电压表和电流表读数都增大}
{电压表和电流表读数都减小}
{电压表读数增大，电流表读数减小}
{电压表读数减小，电流表读数增大}

%% answer: A
%% memo:
\memoanswer{设滑动变阻器的触头上部分电阻为 $x$，则电路的总电阻为 $R_{\text{总}}=r+R_{1}+\frac{xR_{2}}{x+R_{2}}$。滑动变阻器的触头由中点滑向 b 端时，并联支路电阻 $x$ 增大，故路端电压变大，同时并联部分的电压变大，故通过电流表的电流增大，选项 A 正确。}

\item
%% number: 8
%% paperName: 2008年广东理综第 8 题 4 分
%% typeId: 2
%% type: 多选题
%% chapter: 电场
%% point: 电场线
%% method: 无
%% score: 4
%% degree: 650
%% duplicateId: 0
%% body:
图中的实线表示电场线，虚线表示只受电场力作用的带正电粒子的运动轨迹，粒子先经过 M 点，再经过 N 点，可以判定\xzanswer{AD}

\onepicture[w=4.5cm]{figs/08.png}

\fourchoices[answer=AD]
{M 点的电势大于 N 点的电势}
{M 点的电势小于 N 点的电势}
{粒子在 M 点受到的电场力大于在 N 点受到的电场力}
{粒子在 M 点受到的电场力小于在 N 点受到的电场力}

%% answer: AD
%% memo:
\memoanswer{沿着电场线的方向，电势降低，故选项 A 正确。电场线越密，场强越大，同一粒子受到的电场力越大，选项 D 正确。}

\item
%% number: 9
%% paperName: 2008年广东理综第 9 题 4 分
%% typeId: 2
%% type: 多选题
%% chapter: 磁场
%% point: 带电粒子在磁场中的运动
%% method: 无
%% score: 4
%% degree: 650
%% duplicateId: 0
%% body:
带电粒子进入云室会使云室中的气体电离，从而显示其运动轨迹。图是在匀强磁场的云室中观察到的粒子的轨迹，a 和 b 是轨迹上的两点，匀强磁场 $B$ 垂直纸面向里，该粒子在运动时，其质量和电量不变，而动能逐渐减少，下列说法正确的是\xzanswer{AC}

\onepicture[w=3.5cm]{figs/09.png}

\fourchoices[answer=AC]
{粒子先经过 a 点，再经过 b 点}
{粒子先经过 b 点，再经过 a 点}
{粒子带负电}
{粒子带正电}

%% answer: AC
%% memo:
\memoanswer{由 $r=\frac{mv}{qB}$ 可知，粒子的动能越小，圆周运动的半径越小，结合粒子运动轨迹可知，粒子先经过 a 点，再经过 b 点，选项 A 正确。根据左手定则可以判断粒子带负电，选项 C 正确。}

\item
%% number: 10
%% paperName: 2008年广东理综第 10 题 4 分
%% typeId: 2
%% type: 多选题
%% chapter: 直线运动
%% point: 运动图象
%% method: 无
%% score: 4
%% degree: 650
%% duplicateId: 0
%% body:
某人骑自行车在平直道路上行进，图中的实线记录了自行车开始一段时间内的 $v-t$ 图像，某同学为了简化计算，用虚线作近似处理，下列说法正确的是\xzanswer{BD}

\onepicture[w=4.5cm]{figs/10.png}

\fourchoices[answer=BD]
{在 $t_{1}$ 时刻，虚线反映的加速度比实际的大}
{在 $0\sim t_{1}$ 时间内，由虚线计算出的平均速度比实际的大}
{在 $t_{1}\sim t_{2}$ 时间内，由虚线计算出的位移比实际的大}
{在 $t_{3}\sim t_{4}$ 时间内，虚线反映的是匀速运动}

%% answer: BD
%% memo:
\memoanswer{$v-t$ 图线的斜率表示物体的加速度，图线与坐标轴围成的面积表示物体的位移，据此判断选项 B、D 正确。需要注意的是，若为曲线，则曲线切线的斜率表示物体的加速度。}

\item
%% number: 11
%% paperName: 2008年广东理综第 11 题 4 分
%% typeId: 1
%% type: 单选题
%% chapter: 曲线运动
%% point: 平抛运动
%% method: 无
%% score: 4
%% degree: 650
%% duplicateId: 0
%% body:
某同学对着墙壁练习打网球，假定球在墙面上以 $25\Ums$ 速度沿水平方向反弹，落地点到墙面的距离在 $10\Um$ 到 $15\Um$ 之间，忽略空气阻力，球在墙面上反弹点的高度范围是\xzanswer{A}

\fourchoices[answer=A]
{$0.8\Um$ 到 $1.8\Um$}
{$0.8\Um$ 到 $1.6\Um$}
{$1.0\Um$ 到 $1.6\Um$}
{$1.0\Um$ 到 $1.8\Um$}

%% answer: A
%% memo:
\memoanswer{网球反弹后的速度大小几乎不变，故反弹后在空中运动的时间在 $0.4\Us\sim0.6\Us$ 之间。在这个时间范围内，网球下落的高度为 $0.8\Um$ 至 $1.8\Um$；由于竖直方向与地面作用后其速度大小也几乎不变，故还要上升同样的高度，选项 A 正确。}

\item
%% number: 12
%% paperName: 2008年广东理综第 12 题 4 分
%% typeId: 1
%% type: 单选题
%% chapter: 曲线运动
%% point: 变轨问题
%% method: 无
%% score: 4
%% degree: 650
%% duplicateId: 0
%% body:
图是“嫦娥一号奔月”示意图，卫星发射后通过自带的小型火箭多次变轨，进入地月转移轨道，最终被月球引力捕获，成为绕月卫星，并开展对月球的探测。下列说法正确的是\xzanswer{C}

\onepicture[w=6.0cm]{figs/12.png}

\fourchoices[answer=C]
{发射“嫦娥一号”的速度必须达到第三宇宙速度}
{在绕月圆轨道上，卫星周期与卫星质量有关}
{卫星受月球的引力与它到月球中心距离的平方成反比}
{在绕月圆轨道上，卫星受地球的引力大于受月球的引力}

%% answer: C
%% memo:
\memoanswer{由于发射过程中多次变轨，在开始发射时其发射速度必须比第一宇宙速度大，不需要达到第三宇宙速度，选项 A 错误。在绕月轨道上，根据 $F=G\frac{Mm}{r^{2}}=m\frac{4\pi^{2}}{T^{2}}r$ 可知卫星的周期与卫星的质量无关，选项 B 错误，选项 C 正确。由于绕月球运动，地球对卫星的引力较小，故选项 D 错误。}

\item
%% number: 13
%% paperName: 2008年广东理综第 13 题 10 分
%% typeId: 3
%% type: 填空题
%% chapter: 气体性质
%% point: 分子力
%% method: 无
%% score: 10
%% degree: 850
%% duplicateId: 0
%% body:
\begin{enumerate}
	\item
	如图所示，把一块洁净的玻璃板吊在橡皮筋的下端，使玻璃板水平地接触水面，如果你想使玻璃板离开水面，必须用比玻璃板重力 \tkanswer{大} 的拉力向上拉橡皮筋，原因是水分子与玻璃的分子间存在 \tkanswer{引力} 作用。
	\item
	往一杯清水中滴入一滴红墨水，一段时间后，整杯水都变成了红色。这一现象在物理学中称为 \tkanswer{扩散} 现象，是由于分子的 \tkanswer{热运动} 而产生的。这一过程是沿着分子运动的无序性 \tkanswer{增加} 的方向进行的。
\end{enumerate}

\onepicture[align=right]{figs/13.png}

%% answer: 
\jdanswer{
\begin{enumerate}
	\item
	大，引力

	\item
	扩散，热运动，增加

\end{enumerate}
}

%% memo:
\memoanswer{（1）由于水分子和玻璃分子之间存在引力，故需要用比玻璃板重力大的拉力拉橡皮筋，才能使玻璃板离开水面。

（2）由于扩散，红墨水滴入清水后整杯水变红，该过程是由于分子的热运动使无序性增加。}

\item
%% number: 14
%% paperName: 2008年广东理综第 14 题 10 分
%% typeId: 3
%% type: 填空题
%% chapter: 机械振动
%% point: 受迫振动共振
%% method: 无
%% score: 10
%% degree: 850
%% duplicateId: 0
%% body:
\begin{enumerate}
	\item
	大海中航行的轮船，遇到大风大浪冲击时，为了防止倾覆，应当改变航行方向和 \tkanswer{速度}，使风浪冲击力的频率远离轮船摇晃的 \tkanswer{固有频率}。
	\item
	光纤通讯中，光导纤维传递光信号的物理原理是利用光的 \tkanswer{全反射} 现象，要发生这种现象，必须满足的条件是，光从光密介质射向 \tkanswer{光疏介质}，且入射角等于或大于 \tkanswer{临界角}。
\end{enumerate}

%% answer: 
\jdanswer{
\begin{enumerate}
	\item
	速度，固有频率

	\item
	全反射，光疏介质，临界角

\end{enumerate}
}

%% memo:
\memoanswer{（1）为防止轮船发生共振，应改变轮船的航行方向，同时改变轮船的航行速度，使轮船摇摆的频率远离风浪冲击力的频率。

（2）光导纤维利用光的全反射来传递信号。发生全反射时，光必须从光密介质射向光疏介质，且入射角等于或大于临界角。}

\item
%% number: 15
%% paperName: 2008年广东理综第 15 题 11 分
%% typeId: 4
%% type: 实验题
%% chapter: 电路
%% point: 电路实验
%% method: 无
%% score: 11
%% degree: 650
%% duplicateId: 0
%% body:
某实验小组探究一种热敏电阻的温度特性。现有器材：直流恒流电源（在正常工作状态下输出电流恒定）、电压表、待测热敏电阻、保温容器、温度计、开关和导线等。

\begin{enumerate}
	\item
	若用上述器材测量热敏电阻的阻值随温度变化的特性，请在实物图上连线。

	\onepicture[w=6.0cm]{figs/15a.png}

	\item
	实验的主要步骤：

	①正确连接电路，在保温容器中注入适量冷水，接通电源，调节并记录电源的输出电流。

	②在保温容器中添加少量热水，待温度稳定后，闭合开关，\tkanswer{记录电压表电压值}，\tkanswer{记录温度计数值}，断开开关。

	③重复第②步操作若干次，测得多组数据。

	\item
	实验小组算出热敏电阻在不同温度下的阻值，并据此绘得如图所示的 $R-t$ 关系图线，请根据图线写出热敏电阻的 $R-t$ 关系式：$R=$ \tkanswer{100} $+$ \tkanswer{0.395} $t$（$\UO$）（保留三位有效数字）。

	\onepicture[w=8.5cm, align=right]{figs/15b.png}
\end{enumerate}

%% answer: 
\jdanswer{
\begin{enumerate}
	\item
	图略

	\item
	记录电压表电压值、温度计数值

	\item
	$R=100+0.395t$

\end{enumerate}
}

%% memo:
\memoanswer{本题原卷及材料中未提供详解，待补充。}

\item
%% number: 16
%% paperName: 2008年广东理综第 16 题 13 分
%% typeId: 4
%% type: 实验题
%% chapter: 功和能
%% point: 动能定理
%% method: 无
%% score: 13
%% degree: 450
%% duplicateId: 0
%% body:
某实验小组采用图 1 所示的装置探究“动能定理”。图 1 中小车可放置砝码。实验中，小车碰到制动装置时，钩码尚未到达地面。打点计时器工作频率为 $50\UHz$。

\onepicture[w=8.5cm]{figs/16a.png}

\begin{enumerate}
	\item
	实验的部分步骤如下：

	①在小车中放入砝码，把纸带穿过打点计时器，连在小车后端，用细线连接小车和钩码；

	②将小车停在打点计时器附近，\tkanswer{接通电源}，\tkanswer{释放小车}，小车拖动纸带，打点计时器在纸带上打下一列点，\tkanswer{断开电源}；

	③改变钩码或小车中砝码的数量，更换纸带，重复②的操作。

	\item
	图 2 是钩码质量为 $0.03\Ukg$、砝码质量为 $0.02\Ukg$ 时得到的一条纸带，在纸带上选择起始点 O 及 A、B、C、D 和 E 五个计数点，可获得各计数点到 O 的距离 $s$ 及对应时刻小车的速度 $v$，请将 C 点的测量结果填在表 1 中相应的位置。

	\onepicture[w=13.0cm]{figs/16b.png}

	\begin{table}[htbp]
	\centering
	\begin{tblr}{|c|c|c|}
	\hline
	测量点 & $s/\Ucm$ & $v/(\Ums)$ \\
	\hline
	O & 0.00 & 0.35 \\
	\hline
	A & 1.51 & 0.40 \\
	\hline
	B & 3.20 & 0.45 \\
	\hline
	C &  &  \\
	\hline
	D & 7.15 & 0.54 \\
	\hline
	E & 9.41 & 0.60 \\
	\hline
	\end{tblr}
	\end{table}

	\item
	在小车的运动过程中，对于钩码、砝码、小车组成的系统，\tkanswer{钩码的重力}做正功，\tkanswer{小车受摩擦阻力}做负功。
	\item
	实验小组根据实验数据绘出了图 3 中的图线（其中 $\Delta v^{2}=v^{2}-v_{0}^{2}$），根据图线可获得的结论是 \tkanswer{小车初末速度的平方差与位移成正比}，要验证“动能定理”，还需要测量的物理量是摩擦力和 \tkanswer{小车的质量}。

	\onepicture[w=8.0cm, align=right]{figs/16c.png}
\end{enumerate}

%% answer: 
\jdanswer{
\begin{enumerate}
	\item
	②接通电源，释放小车，断开电源

	\item
	5.06，0.49

	\item
	钩码的重力，小车受摩擦阻力

	\item
	小车初末速度的平方差与位移成正比，小车的质量

\end{enumerate}
}

%% memo:
\memoanswer{本题原卷及材料中未提供详解，待补充。}

\item
%% number: 17
%% paperName: 2008年广东理综第 17 题 9 分
%% typeId: 6
%% type: 计算题
%% chapter: 功和能
%% point: 功率
%% method: 无
%% score: 9
%% degree: 650
%% duplicateId: 0
%% body:
为了响应国家的“节能减排”号召，某同学采用了一个家用汽车的节能方法，在符合安全行驶要求的情况下，减少汽车后备箱中放置的不常用物品和控制加油量等措施，使汽车负载减少。假设汽车以 $72\Ukmh$ 的速度匀速行驶时，负载改变前、后汽车所受的阻力分别为 $2000\UN$ 和 $1950\UN$，请计算该方法使汽车发动机输出功率减少了多少？

%% answer: 
\jdanswer{$\Delta P=1\times10^{3}\UW$}

%% memo:
\memoanswer{$v=72\Ukmh=20\Ums$，由 $P=Fv$ 得 $P_{1}=F_{1}v=f_{1}v$，$P_{2}=F_{2}v=f_{2}v$，故 $\Delta P=P_{1}-P_{2}=(f_{1}-f_{2})v=(2000-1950)\times20\UW=1\times10^{3}\UW$。}

\item
%% number: 17
%% paperName: 2008年广东理综第 17 题 9 分
%% typeId: 6
%% type: 计算题
%% chapter: 曲线运动
%% point: 圆周运动的应用
%% method: 无
%% score: 9
%% degree: 650
%% duplicateId: 0
%% body:
有一种叫“飞椅”的游乐项目，示意图如图所示，长为 $L$ 的钢绳一端系着座椅，另一端固定在半径为 $r$ 的水平转盘边缘，转盘可绕穿过其中心的竖直轴转动，当转盘以角速度 $\omega$ 匀速转动时，钢绳与转轴在同一竖直平面内，与竖直方向的夹角为 $\theta$，不计钢绳的重力，求转盘转动的角速度 $\omega$ 与夹角 $\theta$ 的关系。

\onepicture[w=5.0cm, align=right]{figs/17.png}

%% answer: 
\jdanswer{$\omega=\sqrt{\frac{g\tan\theta}{r+L\sin\theta}}$}

%% memo:
\memoanswer{设转盘转动的角速度为 $\omega$ 时，钢绳与竖直方向的夹角为 $\theta$。座椅到中心轴的距离为 $R=r+L\sin\theta$。对座椅分析有 $F_{\text{向}}=mg\tan\theta=mR\omega^{2}=m(r+L\sin\theta)\omega^{2}$，联立两式得 $\omega=\sqrt{\frac{g\tan\theta}{r+L\sin\theta}}$。}

\item
%% number: 18
%% paperName: 2008年广东理综第 18 题 17 分
%% typeId: 6
%% type: 计算题
%% chapter: 电磁感应
%% point: 电磁感应电路分析
%% method: 无
%% score: 17
%% degree: 650
%% duplicateId: 0
%% body:
如图（a）所示，水平放置的两根金属导轨，间距 $L=0.3\Um$，导轨左端连接 $R=0.6\UO$ 的电阻，区域 abcd 内存在垂直于导轨平面 $B=0.6\UT$ 的匀强磁场，磁场区域宽 $D=0.2\Um$，细金属棒 $A_{1}$ 和 $A_{2}$ 用长为 $2D$ 的轻质绝缘杆连接，放置在导轨平面上，并与导轨垂直。每根金属棒在导轨间的电阻为 $r=0.3\UO$，导轨电阻不计，使金属棒以恒定速度 $v=1.0\Ums$ 沿导轨向右穿越磁场。计算从金属棒 $A_{1}$ 进入磁场到 $A_{2}$ 离开磁场的时间内，不同时间段通过电阻 $R$ 的电流强度，并在图（b）中画出。

\onepicture[w=12.0cm, align=right]{figs/18.png}

%% answer: 
\jdanswer{$0\sim t_{1}$（$0\sim0.2\Us$）$I=0.12\UA$；$t_{1}\sim t_{2}$（$0.2\sim0.4\Us$）$I=0$；$t_{2}\sim t_{3}$（$0.4\sim0.6\Us$）$I=0.12\UA$。}

%% memo:
\memoanswer{$A_{1}$ 进入磁场后产生的感应电动势 $E=BDv=0.6\times0.3\times1.0\UV=0.18\UV$。电阻 $R$ 与 $A_{2}$ 并联，阻值 $R_{\text{并}}=\frac{Rr}{R+r}=0.2\UO$，电阻 $R$ 两端电压 $U=\frac{R_{\text{并}}}{R_{\text{并}}+r}E=\frac{0.2}{0.2+0.3}\times0.18\UV=0.072\UV$，通过电阻 $R$ 的电流 $I_{1}=\frac{U}{R}=\frac{0.072}{0.6}\UA=0.12\UA$。

在 $t_{1}\sim t_{2}$（$0.2\sim0.4\Us$）内，$E=0$，$I_{2}=0$。

在 $t_{2}\sim t_{3}$（$0.4\sim0.6\Us$）内，同理 $I_{3}=0.12\UA$。}

\item
%% number: 19
%% paperName: 2008年广东理综第 19 题 16 分
%% typeId: 6
%% type: 计算题
%% chapter: 运动和力
%% point: 动量守恒定律
%% method: 无
%% score: 16
%% degree: 450
%% duplicateId: 0
%% body:
如图（a）所示，在光滑绝缘水平面的 AB 区域内存在水平向右的电场，电场强度 $E$ 随时间的变化如图（b）所示。不带电的绝缘小球 $P_{2}$ 静止在 O 点，$t=0$ 时，带正电的小球 $P_{1}$ 以速度 $v_{0}$ 从 A 点进入 AB 区域，随后与 $P_{2}$ 发生正碰后反弹，反弹速度大小是碰前的 $\frac{2}{3}$ 倍，$P_{1}$ 的质量为 $m_{1}$、带电量为 $q$，$P_{2}$ 的质量为 $m_{2}=5m_{1}$。AO 间距为 $L_{0}$，OB 间距 $L=\frac{4}{3}L_{0}$，已知 $\frac{qE_{0}}{m_{1}}=\frac{2v_{0}^{2}}{3L_{0}}$，$T=\frac{L_{0}}{v_{0}}$。

\begin{enumerate}
	\item
	求碰撞后小球 $P_{1}$ 向左运动的最大距离及所需时间；
	\item
	讨论两球是否能在 OB 区间内再次发生碰撞。
\end{enumerate}

\onepicture[w=13.0cm, align=right]{figs/19.png}

%% answer: 
\jdanswer{
\begin{enumerate}
	\item
	$s_{m}=\frac{L_{0}}{3}$，$t_{2}=\frac{L_{0}}{v_{0}}$

	\item
	两球能在 OB 区间内再次碰撞

\end{enumerate}
}

%% memo:
\memoanswer{（1）$P_{1}$ 经 $t_{1}=\frac{L_{0}}{v_{0}}$ 时间与 $P_{2}$ 碰撞。设碰后 $P_{2}$ 的速度为 $v_{2}$，由动量守恒 $m_{1}v_{0}=m_{1}\left(-\frac{2}{3}v_{0}\right)+m_{2}v_{2}$，解得 $v_{2}=\frac{1}{3}v_{0}$，方向水平向右。碰后 $P_{1}$ 向左运动，加速度 $a_{1}=\frac{qE_{0}}{m_{1}}=\frac{2v_{0}^{2}}{3L_{0}}$，最大距离 $s_{m}=\frac{v_{1}^{2}}{2a_{1}}=\frac{L_{0}}{3}$，所用时间 $t_{2}=\frac{v_{1}}{a_{1}}=\frac{L_{0}}{v_{0}}$。

（2）设 $P_{1}$、$P_{2}$ 碰撞后又经 $\Delta t$ 时间在 OB 区间内再次发生碰撞，且 $P_{1}$ 受电场力不变，以水平向右为正，则 $-v_{1}\Delta t+\frac{1}{2}a_{1}\Delta t^{2}=v_{2}\Delta t$，解得 $\Delta t=\frac{3L_{0}}{v_{0}}=3T$，可见电场力不变。对 $P_{2}$ 分析：$s_{2}=v_{2}\Delta t=\frac{1}{3}v_{0}\times\frac{3L_{0}}{v_{0}}=L_{0}<L=\frac{4}{3}L_{0}$，所以假设成立，两球能在 OB 区间内再次碰撞。}

\item
%% number: 20
%% paperName: 2008年广东理综第 20 题 18 分
%% typeId: 6
%% type: 计算题
%% chapter: 运动和力
%% point: 动量守恒定律
%% method: 无
%% score: 18
%% degree: 450
%% duplicateId: 0
%% body:
如图所示，固定的凹槽水平表面光滑，其内放置 U 形滑板 N，滑板两端为半径 $R=0.45\Um$ 的四分之一圆弧面，A 和 D 分别为圆弧的端点，BC 段表面粗糙，其余段表面光滑，小滑块 $P_{1}$ 和 $P_{2}$ 的质量均为 $m$，滑板的质量 $M=4m$，$P_{1}$ 和 $P_{2}$ 与 BC 面的动摩擦因数分别为 $\mu_{1}=0.10$ 和 $\mu_{2}=0.40$，最大静摩擦力近似等于滑动摩擦力。开始时滑板紧靠槽的左端，$P_{2}$ 静止在粗糙面的 B 点，$P_{1}$ 以 $v_{0}=4.0\Ums$ 的初速度从 A 点沿弧面自由滑下，与 $P_{2}$ 发生弹性碰撞后，$P_{1}$ 处在粗糙面的 B 点上，当 $P_{2}$ 滑到 C 点时，滑板恰好与槽的右端碰撞并与槽牢固粘连，$P_{2}$ 继续滑动，到达 D 点时速度为零，$P_{1}$ 与 $P_{2}$ 视为质点，取 $g=10\Umsq$。问：

\begin{enumerate}
	\item
	$P_{2}$ 在 BC 段向右滑动时，滑板的加速度为多大？
	\item
	BC 长度为多少？N、$P_{1}$ 和 $P_{2}$ 最终静止后，$P_{1}$ 和 $P_{2}$ 间的距离为多少？
\end{enumerate}

\onepicture[w=13.0cm, align=right]{figs/20.png}

%% answer: 
\jdanswer{
\begin{enumerate}
	\item
	$a_{2}=0.8\Umsq$

	\item
	$L=1.9\Um$，$\Delta s=0.695\Um$

\end{enumerate}
}

%% memo:
\memoanswer{（1）$P_{1}$ 滑到最低点时速度为 $v_{1}$，由机械能守恒 $\frac{1}{2}mv_{0}^{2}+mgR=\frac{1}{2}mv_{1}^{2}$，解得 $v_{1}=5\Ums$。$P_{1}$、$P_{2}$ 发生弹性碰撞，由动量守恒和机械能守恒解得 $v_{1}'=0$，$v_{2}'=5\Ums$，$P_{2}$ 向右运动。假设 $P_{1}$ 不动，对 $P_{2}$ 有 $f_{2}=\mu_{2}mg=4m$，对 $P_{1}$ 和 N 有 $f_{2}=(m+M)a_{2}$，得 $a_{2}=0.8\Umsq$。此时 $P_{1}$ 所需的摩擦力 $f=ma_{2}=0.8m<f_{m}=1.0m$，所以假设成立。

（2）$P_{2}$ 滑到 C 点时速率为 $v_{2}'$，由 $mgR=\frac{1}{2}mv_{2}'^{2}$ 得 $v_{2}'=3\Ums$。设此时 $P_{1}$ 和 N 的速率为 $v$，由动量守恒 $mv_{2}=(m+M)v+mv_{2}'$，得 $v=0.4\Ums$。对 $P_{1}$、$P_{2}$、M 组成的系统，由能量守恒 $f_{2}L=\frac{1}{2}mv_{2}^{2}-\frac{1}{2}mv_{2}'^{2}-\frac{1}{2}(m+M)v^{2}$，解得 $L=1.9\Um$。滑板碰后，$P_{1}$ 向右滑行的距离 $s_{1}=\frac{v^{2}}{2a_{1}}=0.08\Um$，$P_{2}$ 向左滑行的距离 $s_{2}=\frac{v_{2}'^{2}}{2a_{2}}=1.125\Um$，所以 $P_{1}$、$P_{2}$ 静止后的距离 $\Delta s=L-s_{1}-s_{2}=0.695\Um$。}

\end{enumerate}

\end{document}
